
\section{The Cross Product (Vector Product)}
The dot product measures alignment between vectors. Many physical problems instead require a vector perpendicular to a plane, an oriented area, or a rotational effect. This is provided by the \emph{cross product}.

\subsection{Geometric Definition}
For vectors $\mathbf A$ and $\mathbf B$ separated by an angle $\theta$,
\[
\boxed{\mathbf A\times\mathbf B=|\mathbf A||\mathbf B|\sin\theta\,\hat{\mathbf n}}
\]
where $\hat{\mathbf n}$ is a unit vector perpendicular to the plane containing $\mathbf A$ and $\mathbf B$.

The magnitude equals the area of the parallelogram spanned by the two vectors,
\[
|\mathbf A\times\mathbf B|=|\mathbf A||\mathbf B|\sin\theta.
\]

\subsection{Right-Hand Rule}
The direction is determined by the right-hand rule.

\subsection{Basis Relations}
\[
\mathbf e_x\times\mathbf e_y=\mathbf e_z,\qquad
\mathbf e_y\times\mathbf e_z=\mathbf e_x,\qquad
\mathbf e_z\times\mathbf e_x=\mathbf e_y.
\]
Reversing the order changes the sign,
\[
\mathbf A\times\mathbf B=-\mathbf B\times\mathbf A.
\]

\subsection{Component Formula}
For
\[
\mathbf A=(A_x,A_y,A_z),\qquad
\mathbf B=(B_x,B_y,B_z),
\]
\[
\mathbf A\times\mathbf B=
\begin{pmatrix}
A_yB_z-A_zB_y\\
A_zB_x-A_xB_z\\
A_xB_y-A_yB_x
\end{pmatrix}.
\]

\subsection{Applications}
\[
\boldsymbol\tau=\mathbf r\times\mathbf F,\qquad
\mathbf L=\mathbf r\times\mathbf p,\qquad
\mathbf F=q\,\mathbf v\times\mathbf B,\qquad
\mathbf S=\mathbf E\times\mathbf H.
\]

\subsection{Worked Example}
For $\mathbf A=(2,1,0)$ and $\mathbf B=(1,3,4)$,
\[
\mathbf A\times\mathbf B=(4,-8,5),
\]
which satisfies $(\mathbf A\times\mathbf B)\cdot\mathbf A=0$ and
$(\mathbf A\times\mathbf B)\cdot\mathbf B=0$.

\subsection{Looking Ahead}
The next section develops vector projections and orthogonal decomposition.
